% Source: OpenAI2026AreaLaw, 06-transport.tex, transport:states, lines 393--401.
% Manuscript revision: adc7f1241b42e322a6451854ab7e4b4c146bf78a.
\section{Embedding identities for trace-adjoint leaf maps}
\label{sec:leaf_state_intertwine}

Let $T$ be a weighted binary tree with labels in a set $I$.
For positive definite inputs $(A_i)_{i\in I}$, let $\Phi_{T,j}^A$
be its normalized leaf map, as in
Theorem~\ref{thm:mean_tree_derivative_intertwine}.
The trace-pairing adjoint is characterized by
\begin{align}
  \operatorname{Tr}\bigl((\Phi_{T,j}^A)^*(\rho)X\bigr)
    &=\operatorname{Tr}\bigl(\rho\,\Phi_{T,j}^A(X)\bigr).
\end{align}
The following identity is auxiliary to the transport-state construction
of~\cite{OpenAI2026AreaLaw}, display \texttt{transport:states}.
It applies to arbitrary matrices; positivity and trace normalization
may be imposed separately when constructing states.

\begin{theorem}[Embedding a matrix through a trace-adjoint leaf map]
  \label{thm:trace_adjoint_leaf_map_intertwine}
  \lean{Matrix.MeanTree.traceAdjoint_leafMap_intertwine}
  \leanok
  Let $(A_i)_{i\in I}$ and $(B_i)_{i\in I}$ be positive definite inputs
  on two finite-dimensional complex Hilbert spaces.
  Suppose that a linear map $J$ from the second space to the first
  satisfies $A_iJ=JB_i$ for every $i\in I$.
  Then, for the same tree $T$, every label $j$ and every matrix $\rho$
  on the second space satisfy
  \begin{align}
    (\Phi_{T,j}^A)^*(J\rho J^*)
      &=J(\Phi_{T,j}^B)^*(\rho)J^*.
  \end{align}
  The inputs use the same labels and weights, including repeated labels
  and zero terminal weights. No isometry assumption on $J$ is required.
\end{theorem}

\begin{proof}
  \leanok
  \uses{thm:mean_tree_derivative_compress}
  Put $C_J(X)=J^*XJ$ and $E_J(\rho)=J\rho J^*$.
  Cyclicity of the trace gives $C_J^*=E_J$.
  Theorem~\ref{thm:mean_tree_derivative_compress} gives
  \begin{align}
    C_J\circ\Phi_{T,j}^A&=\Phi_{T,j}^B\circ C_J.
  \end{align}
  The trace-pairing adjoint reverses composition. Taking adjoints
  therefore gives
  \begin{align}
    (\Phi_{T,j}^A)^*\circ E_J
      &=E_J\circ(\Phi_{T,j}^B)^*.
  \end{align}
  Evaluate this identity at $\rho$.
\end{proof}
